\begin{figure}[!b]
\centering
\begin{tikzpicture}[x=1cm,y=1cm,>=Stealth,font=\small]
  \fill[blue!7] (0,0) rectangle (2.3,2.2);
  \fill[orange!12] (2.3,0) rectangle (5,2.2);
  \draw[thick] (0,0) -- (5,0);
  \draw[thick] (0,2.2) -- (5,2.2);
  \draw[densely dashed,thick] (2.3,0) -- (2.3,2.2);
  \node[align=center] at (1.1,1.2) {retained\\closed side};
  \node[align=center] at (3.6,1.2) {$L\times[0,\infty)$\\attached\\cylinder};
  \draw[->,thick] (4.4,.5) -- (5.3,.5) node[right] {$s$};
  \node[align=center,below] at (2.3,-.1) {internal attachment seam};
  \draw[very thick,teal!70!black] (7,0) -- (7,2.2);
  \node[align=center,right] at (7.2,1.1) {separate hit space $L$\\base dimension $m-1$};
  \node[above] at (2.5,2.7) {either sign side: base dimension $m$};
\end{tikzpicture}
\caption{A regular cut retains each closed sign side with a full cylinder
attached, and also retains the wall as a separate hit space of one lower base
dimension. The attachment seam is internal to the extended side. In old
product ends this entire operation is a product. The raw metric and its
later smoothing are defined in Appendices~\ref{sec:cut} and~\ref{sec:regularization}.}
\label{fig:regular-packet-split}
\end{figure}
